\begin{table}[!htbp]
\centering
\caption[Regresie numerică cross-limbaj]{Verificare numerică de regresie cross-limbaj: abaterea maximă absolută între traiectoriile compartimentale produse în R, Julia și Python pentru aceiași parametri și aceleași condiții initiale, pe setul sintetic \texttt{multiwave} (populația $N=50 \cdot 10^6$, $t_{\max}=365$ pași zilnici). Modelele folosesc solvere diferite (R cu RK4 pas fix $h=1$, Python cu RK45 adaptiv, Julia cu Tsit5), deci abaterile reflectă toleranțele numerice cumulate, nu erori de implementare.}
\label{tab:cross_language_regression}
\resizebox{\columnwidth}{!}{%
\begin{tabular}{l l r r}
\toprule
\textbf{Model} & \textbf{Pereche} & \textbf{Max $|\Delta|$} & \textbf{Normalizat ($/N$)} \\
\midrule
\texttt{sir} & \texttt{julia} vs \texttt{python} & 0.5 & $1{,}01 \cdot 10^{-8}$ \\
 & \texttt{julia} vs \texttt{r} & 2185.8 & $4{,}37 \cdot 10^{-5}$ \\
 & \texttt{python} vs \texttt{r} & 2186.2 & $4{,}37 \cdot 10^{-5}$ \\
\midrule
\texttt{sidarthe} & \texttt{julia} vs \texttt{python} & 0.2 & $3{,}72 \cdot 10^{-9}$ \\
 & \texttt{julia} vs \texttt{r} & 6570.9 & $1{,}31 \cdot 10^{-4}$ \\
 & \texttt{python} vs \texttt{r} & 6571.1 & $1{,}31 \cdot 10^{-4}$ \\
\bottomrule
\end{tabular}%
}
\end{table}
